\chapter{总结与展望}
\section{工作总结}

随着模型规模与训练并行度的持续提升，通信时间在端到端训练时间中的占比不断上升，逐渐成为影响分布式训练效率的主要瓶颈。围绕这一瓶颈，本文从环形集合通信算法优化设计、拓扑感知的集合通信算法自动合成以及张量并行中的计算通信并行方法三个层面展开研究，致力于在真实训练环境中获得稳定且可扩展的性能提升。

本文的主要工作总结如下：

（1）提出了链路聚合的双半环集合通信算法，并完成工程实现与验证。该算法的核心思想是在双向链路可并行传输的环形互连环境中,通过重构双向环算法的数据流传输顺序使两组反向链路协同工作：在主体通信阶段传输完整数据块，仅在最后一步对对角节点数据进行切分，从而将通信轮数从 $n-1$ 压缩至约 $n/2$，同时将单次传输的数据包规模扩大一倍以提升链路带宽利用率。本文给出了 AllGather、ReduceScatter 与 AllReduce 三类算子的统一实现方式，并分析了其在静态时延与带宽利用方面的理论优势。在华为昇腾 910C 集群上基于 HCCL 完成实现与评估，结果表明双半环算法在 16/32/64 节点、64\,MB 至 2\,GB 数据量配置下均稳定优于基线的双向环算法，实现 2.7\% 至 33.1\% 的算法带宽提升；在大规模与大数据量场景中收益较小，但整体性能始终优于基线方案。

（2）提出了交换机拓扑感知的通用集合通信算法合成方法 FTACAS。该方法以时间扩展网络（TEN）为建模基础，将物理拓扑沿时间维度展开，统一刻画设备节点、交换机节点、链路带宽与静态时延，并显式建模 Cut-Through 与 Store-and-Forward 两种转发模式下交换机的端口并发与缓冲约束。在此之上，FTACAS 采用事件驱动的贪心合成框架：以“数据块到达设备节点”为事件触发点，动态构建可用链路的匹配空间，并按贪心策略选择下一跳链路，从而在多项式时间内生成可行且高质量的通信算法。针对 AllGather 与 AlltoAll 两类算子，本文分别给出了特定的匹配策略。在多种典型 GPU 集群拓扑上的对比实验表明：在 AllGather 场景下，FTACAS 的算法带宽整体不低于 TECCL，在部分复杂拓扑上可达到其 2 至 3 倍；在 AlltoAll 场景下，FTACAS 在多数拓扑中与 TECCL 接近，最低也能达到 TECCL 的 89\% 以上，最高可达到1.5倍。此外，FTACAS 的合成时间相比 TECCL 普遍降低数百至数千倍，最高可达约 $7.79\times10^{4}$ 倍，显著增强了算法合成在工程实践中的可用性。

（3）提出了面向张量并行训练的计算通信并行优化方法，并在 TP/DP 混合并行下验证其端到端收益。针对 Megatron 张量并行方法中多处同步通信位于关键路径、易导致设备空闲的问题，本文通过将 AllReduce 拆分为 ReduceScatter 与 AllGather，并结合矩阵乘法的可分块特性，在前向与反向传播阶段分别设计了细粒度调度策略。在前向传播中，通过重排矩阵乘法的分块计算顺序，使行并行阶段的计算与 ReduceScatter 通信实现流水式重叠，并在列并行阶段利用 AllGather 接收数据的时序特性同步启动对应分块计算；在反向传播中，进一步利用输入梯度与权重梯度计算的独立性，将通信操作与梯度计算并行执行，减少同步等待时间。基于 PyTorch 接口实现后，在单节点 8 卡 NVIDIA L40S 环境下以 TP=4、DP=2 的混合并行配置进行实测，三种不同模型规格的端到端训练时间分别获得 1.25$\times$、1.29$\times$、1.25$\times$ 加速，约 45\% 至 53\% 的 TP 侧通信开销被成功从关键路径中隐藏，端到端训练时间显著缩短。



\section{研究展望}

本文围绕集合通信与通算协同给出了若干优化方法，但面向更大规模、更复杂互连以及更丰富并行策略的训练系统，仍有进一步拓展空间。结合分布式训练的发展趋势，未来可从以下几个方面开展深入研究：

（1）面向更复杂互连环境的集合通信算法扩展与自适应选择。双半环算法在双向并行链路的环形互连中取得了稳定收益，后续可进一步研究其在多级互连与链路异构场景下的分层扩展策略，例如将机内高带宽域与机间网络分离建模，在不同层级采用差异化的数据切分与传输节奏；同时结合消息大小与链路状态设计自适应的算法/分块粒度选择机制，使通信库能够在运行时根据代价模型自动切换最合适的实现，从而提升方法在常见硬件平台上的鲁棒性与通用性。

（2）集合通信算法合成的模型精化与算子覆盖扩展。FTACAS 当前以 TEN 与事件驱动的贪心合成框架为核心，在合成效率上具备明显优势，但在高复杂度算子（如大规模 AlltoAll）或极端拓扑下仍可能受到局部最优影响。后续可考虑引入更精细的代价模型与校准机制，将链路拥塞、交换机端口竞争、多链路并发等因素纳入建模；在求解层面，可探索先以启发式快速生成可行算法、再结合局部改进或混合求解提升质量的组合策略，兼顾算法质量与生成时间。此外，进一步扩展对 AllReduce、ReduceScatter、Broadcast 等更多基础算子的统一支持，并与实际通信库的实现约束（如消息分块粒度、缓冲区限制、内存对齐）对接，将有助于提升合成器在训练系统中的适用范围。

（3）张量并行通算协同的自动化调度与跨节点验证。本文的计算通信并行方法在单节点 TP/DP 混合并行下验证了收益，未来可面向多节点环境进一步评估与优化，重点关注机间网络时延、带宽波动以及拓扑差异对重叠窗口的影响；同时，结合算子融合与编译优化，减少拆分带来的额外算子调用开销与算子效率下降问题，并探索自适应分块以在不同批次大小/序列长度下保持稳定加速。更进一步，可将计算通信调度与拓扑感知的集合通信算法（如 FTACAS 合成的算法）相结合，使计算与合成的定制通信算法协同执行，形成端到端优化闭环，以支撑更高维度的混合并行与更大规模的训练任务。
